\Article{20}
\textsc{Aufgabe.} \textit{Es bezeichnen \(v,v',v''\) etc. die folgenden linearen Funktionen der Unbestimmten \(x,y,z\) etc.}
\[
\left.\begin{aligned}
v&=ax+by+cz+\text{etc.}+l\\
v'&=a'x+b'y+c'z+\text{etc.}+l'\\
v''&=a''x+b''y+c''z+\text{etc.}+l''\text{ etc.}
\end{aligned}\right\}\quad\text{(I)}
\]
\textit{Aus allen Systemen von Koeffizienten \(\kappa,\kappa',\kappa''\) etc., für die identisch}
\[
\kappa v+\kappa'v'+\kappa''v''+\text{etc.}=x-k
\]
\textit{gilt, wobei \(k\) eine bestimmte, das heißt von \(x,y,z\) etc. unabhängige Größe ist, ist dasjenige zu ermitteln, für das \(\kappa\kappa+\kappa'\kappa'+\kappa''\kappa''+\text{etc.}\) den kleinsten Wert annimmt.}

\textit{Lösung.} Setzen wir
\[
\left.\begin{aligned}
av+a'v'+a''v''+\text{etc.}&=\xi\\
bv+b'v'+b''v''+\text{etc.}&=\eta\\
cv+c'v'+c''v''+\text{etc.}&=\zeta
\end{aligned}\right\}\quad\text{(II)}
\]
etc.: Dann sind auch \(\xi,\eta,\zeta\) etc. lineare Funktionen von \(x,y,z\) etc., nämlich
\[
\left.\begin{aligned}
\xi&=x\Sigma aa+y\Sigma ab+z\Sigma ac+\text{etc.}+\Sigma al\\
\eta&=x\Sigma ab+y\Sigma bb+z\Sigma bc+\text{etc.}+\Sigma bl\\
\zeta&=x\Sigma ac+y\Sigma bc+z\Sigma cc+\text{etc.}+\Sigma cl\text{ etc.}
\end{aligned}\right\}\quad\text{(III)}
\]
(wobei \(\Sigma aa\) die Summe \(aa+a'a'+a''a''+\text{etc.}\) bezeichnet und Entsprechendes für die übrigen gilt), und die Anzahl von \(\xi,\eta,\zeta\) etc. ist der Anzahl der Unbestimmten \(x,y,z\) etc. gleich, nämlich \(=\rho\). Durch Elimination lässt sich daher eine Gleichung der folgenden Art gewinnen\footnote{Warum wir gerade diese Zeichen zur Bezeichnung der aus einer solchen Elimination hervorgehenden Koeffizienten gewählt haben, wird unten klar werden.}
\SourceLeaf{028}
\[
x=A+[\alpha\alpha]\xi+[\alpha\beta]\eta+[\alpha\gamma]\zeta+\text{etc.}
\]
in der eine identische Gleichung entstehen muss, wenn für \(\xi,\eta,\zeta\) etc. ihre Werte aus III eingesetzt werden. Setzen wir also
\[
\left.\begin{aligned}
a[\alpha\alpha]+b[\alpha\beta]+c[\alpha\gamma]+\text{etc.}&=\alpha\\
a'[\alpha\alpha]+b'[\alpha\beta]+c'[\alpha\gamma]+\text{etc.}&=\alpha'\\
a''[\alpha\alpha]+b''[\alpha\beta]+c''[\alpha\gamma]+\text{etc.}&=\alpha''\text{ etc.}
\end{aligned}\right\}\quad\text{(IV)}
\]
so gilt notwendig identisch
\[
\alpha v+\alpha'v'+\alpha''v''+\text{etc.}=x-A\quad\text{(V)}
\]
Diese Gleichung zeigt, dass zu den Systemen von Werten der Koeffizienten \(\kappa,\kappa',\kappa''\) etc. gewiss auch \(\kappa=\alpha,\kappa'=\alpha',\kappa''=\alpha''\) etc. gehören und dass für jedes beliebige System identisch
\[
(\kappa-\alpha)v+(\kappa'-\alpha')v'+(\kappa''-\alpha'')v''+\text{etc.}=A-k
\]
gelten muss. Diese Gleichung schließt die folgenden ein
\[
\begin{aligned}
(\kappa-\alpha)a+(\kappa'-\alpha')a'+(\kappa''-\alpha'')a''+\text{etc.}&=0\\
(\kappa-\alpha)b+(\kappa'-\alpha')b'+(\kappa''-\alpha'')b''+\text{etc.}&=0\\
(\kappa-\alpha)c+(\kappa'-\alpha')c'+(\kappa''-\alpha'')c''+\text{etc.}&=0\text{ etc.}
\end{aligned}
\]
Multiplizieren wir diese Gleichungen jeweils mit \([\alpha\alpha],[\alpha\beta],[\alpha\gamma]\) etc. und addieren, so erhalten wir wegen (IV)
\[
(\kappa-\alpha)\alpha+(\kappa'-\alpha')\alpha'+(\kappa''-\alpha'')\alpha''+\text{etc.}=0
\]
oder, was dasselbe ist,
\[
\begin{aligned}
&\kappa\kappa+\kappa'\kappa'+\kappa''\kappa''+\text{etc.}\\
&=\alpha\alpha+\alpha'\alpha'+\alpha''\alpha''+\text{etc.}{}\\
&\quad+(\kappa-\alpha)^2+(\kappa'-\alpha')^2+(\kappa''-\alpha'')^2+\text{etc.}
\end{aligned}
\]
woraus ersichtlich ist, dass die Summe \(\kappa\kappa+\kappa'\kappa'+\kappa''\kappa''+\text{etc.}\) ihren kleinsten Wert annimmt, wenn \(\kappa=\alpha,\kappa'=\alpha',\kappa''=\alpha''\) etc. gesetzt wird. Q.~E.~I.

Dieser kleinste Wert selbst lässt sich im Übrigen auf folgende Weise ermitteln. Aus Gleichung (V) folgt
\SourceLeaf{029}
\[
\begin{aligned}
\alpha a+\alpha'a'+\alpha''a''+\text{etc.}&=1\\
\alpha b+\alpha'b'+\alpha''b''+\text{etc.}&=0\\
\alpha c+\alpha'c'+\alpha''c''+\text{etc.}&=0\text{ etc.}
\end{aligned}
\]
Multiplizieren wir diese Gleichungen jeweils mit \([\alpha\alpha],[\alpha\beta],[\alpha\gamma]\) etc. und addieren, so erhalten wir mithilfe der Gleichungen (IV) unmittelbar
\[
\alpha\alpha+\alpha'\alpha'+\alpha''\alpha''+\text{etc.}=[\alpha\alpha]
\]

\clearpage
